\section{Trace all five rows}
\Needspace{10\baselineskip}
\begin{longtable}{p{11mm}p{17mm}p{54mm}p{65mm}}
  \toprule Row & Filter  & Inflated sample residuals (decimal) & Reconstructed samples (decimal) \\\midrule
  0            & None    & 10 20 30 40 50 60                   & 10 20 30 40 50 60               \\
  1            & Up      & 5 5 5 5 5 5                         & 15 25 35 45 55 65               \\
  2            & Sub     & 20 30 40 30 30 30                   & 20 30 40 50 60 70               \\
  3            & Average & 15 20 25 18 18 18                   & 25 35 45 55 65 75               \\
  4            & Paeth   & 5 5 5 5 5 5                         & 30 40 50 60 70 80               \\\bottomrule
\end{longtable}
For row two, the first red residual is 20 and has no left neighbor, so reconstructed
red is 20. The second red residual is 30. Its left neighbor is now the reconstructed
20, giving 50. Do not add the preceding blue byte, and do not use an unprocessed
residual from the left pixel.

For the second red sample of row three, $a=25$ and $b=50$. The predictor is
$\lfloor75/2\rfloor=37$, and residual 18 restores 55. This also explains why reading
the whole row into vector lanes and adding a vector of left residuals is incorrect:
the left values needed for prediction may only become available during reconstruction.

The compact RGB bytes for the first pixel are hexadecimal \texttt{0A 14 1E}.
Its display word is \texttt{0x000A141E}. In little-endian memory that word occupies
\texttt{1E 14 0A 00}; in big-endian memory it occupies \texttt{00 0A 14 1E}.
The word's channel interpretation is unchanged. The test must compare the numerical
framebuffer words, with byte-order-specific byte checks when examining raw memory.
